% ====================================================================
% s34_mech.tex — §3.4 기계 히스테리시스(GS-1): Larché–Cahn 시도 (W2 초안 — 실측·데이터 강조축)
%   유도 시도: 닫히는 데까지만(고정 응력 강체 중심이동). 못 닫으면 공백 4분류 선언(유사 유도 날조 금지).
%   강조: 실측 앵커(100–120 mV/GPa·−1.75 GPa)로 닫힘 지점·미닫힘 지점을 데이터로 국소화.
% ====================================================================
\section{기계 히스테리시스 --- 응력--화학퍼텐셜 결합의 유도 시도 (GS-1)}\label{sec:si-mech}
생존 지도 N3 는 Si 히스테리시스를 유일한 ``새 물리'' 로 분류했다: 정규용액 spinodal gap 상한($\Omega{=}4RT$
에서 $\approx55$ mV, 그림~\ref{fig:hysgap})으로는 실측 수백 mV 를 담을 수 없고, in~situ 측정이 보이는 것은
자유에너지 이중웰이 아니라 소성 유동과 응력--전위 결합이기 때문이다\cite{sethuraman_stressevo2010,
sethuraman_stresspot2010}. 골격에 없는 항은 \emph{응력이 화학퍼텐셜에 들어가는 항}이다. 본 절은 그 항의
표준 이론틀(Larché--Cahn)을 본문 결과 사슬에 끌어와 \emph{닫히는 데까지} 유도하고, 닫히지 않는 지점을
날조 없이 공백으로 선언한다.

\subsection{Larché--Cahn 선형 결합 --- 닫히는 부분}\label{ssec:si-lc}
\textbf{(a) 출발.} 개방계 고체에서 삽입종의 화학퍼텐셜은 응력에 선형으로 결합한다. Larché--Cahn 선형
이론\cite{larchecahn1973}의 표준 출발형은, host 내 Li 부분몰 부피 $V_\mathrm{Li}$ 와 정수압(평균) 응력
$\sigma_h=\tfrac13\operatorname{tr}\sigma$ 로
\begin{equation}
\mu_\mathrm{Li}(\theta,\sigma)\;=\;\mu_\mathrm{Li}(\theta,0)\;-\;V_\mathrm{Li}\,\sigma_h
\label{eq:si-lc-mu}
\end{equation}
--- 압축($\sigma_h<0$)이면 삽입종 화학퍼텐셜이 오르고(리튬화가 억제되고), 인장이면 내린다는 부호다.
\textbf{(b) 결선.} 이 $\mu$ 를 본문 전기화학 평형 조건 식~\eqref{eq:sm-eqcond}($\mu_\mathrm{Li}(\theta_\eq)=\mu^0-sF(V-U_j)$)에
넣으면 응력 항이 상수 덩이에 얹혀
\begin{equation}
\mu_\mathrm{Li}(\theta_\eq,0)\;=\;\mu^0-sF\bigl(V-U_j-\Delta U_\sigma\bigr),
\qquad
\Delta U_\sigma\equiv\frac{V_\mathrm{Li}\,\sigma_h}{sF}
\label{eq:si-lc-shift}
\end{equation}
가 된다. \textbf{(c) 닫힘.} 곧 \emph{$\sigma_h$ 가 고정 상수라면} 응력은 평형 중심을 $\Delta U_\sigma$ 만큼
강체 이동시킬 뿐이고, 평형 종$\cdot$logistic 은 중심만 옮긴 채 그대로 닫힌다 --- 식~\eqref{eq:sm-nernst}
꼴에서 $U_j\to U_j+\Delta U_\sigma$. 여기까지가 1--2 스텝에 닫히는 부분이다.
\begin{srcbox}
실측 앵커(닫힘 지점) --- $\Delta U_\sigma$ 의 기울기 $\partial V/\partial\sigma_h=V_\mathrm{Li}/F$ 는
in~situ 측정 응력--전위 결합 $\sim100$--$120$ mV/GPa\cite{sethuraman_stresspot2010}(tier A)와 같은
규모다(부분몰 부피 스케일의 실측 앵커). 단 $V_\mathrm{Li}$ 의 절대 수치는 본 조사 표에 없어 여기서
고정하지 않는다(확인 필요) --- 기울기의 규모 정합까지만 tier-A 로 읽는다.
\end{srcbox}

\subsection{정직 공백 GS-1 --- 닫히지 않는 부분}\label{ssec:si-gs1}
히스테리시스는 중심의 \emph{단일 이동}이 아니라 충$\cdot$방전 사이의 \emph{갈림}이다. 식~\eqref{eq:si-lc-shift}
이 갈림을 내려면 $\sigma_h$ 가 리튬화(압축)와 탈리튬화(인장)에서 \emph{다른 경로}를 밟아야 하고, 그 경로
차이가 곧 기계 히스 gap 이다:
\begin{equation}
\Delta U_\hys^{\mathrm{mech}}(\theta)\;\approx\;\frac{V_\mathrm{Li}}{sF}\Bigl[\sigma_h^{\chg}(\theta)-\sigma_h^{\dis}(\theta)\Bigr].
\label{eq:si-lc-gap}
\end{equation}
문제는 $\sigma_h^{\chg}\ne\sigma_h^{\dis}$ 를 낳는 것이 \emph{소성 유동}이라는 데 있다 --- 실측은 리튬화 시
압축 소성 유동이 $\sim-1.75$ GPa 에 이르고\cite{sethuraman_stressevo2010}(tier A), 이 항복$\cdot$이력
의존성이 응력을 경로$\cdot$이력의 함수 $\sigma_h(\theta,\text{이력})$ 로 만든다. 이 구성식(항복 응력$\cdot$유동
법칙의 탄소성 모델)은 평형 이중웰 열역학의 밖이며, 식~\eqref{eq:si-lc-gap} 을 닫으려면 그 구성식이 있어야
한다. 최근 모델은 이를 \emph{수치로} 닫는다: 입자--SEI 점탄소성 chemo-mechanical
모델\cite{koebbing2024} 과 다단 상전이$\cdot$(비)결정화 경로분기 현상론\cite{jiang_sihys2020} 이 각각
갈림을 재현한다.

\begin{keybox}
\textbf{GS-1 공백 4분류.} 닫힌 부분: 고정 응력의 \emph{강체 중심 이동} $\Delta U_\sigma$(식~\eqref{eq:si-lc-shift})
와 히스 gap 의 \emph{구조}(경로차, 식~\eqref{eq:si-lc-gap}) --- Larché--Cahn 선형 결합\cite{larchecahn1973}
으로 1--2 스텝 닫힘. 미닫힘 부분의 분류는 둘이 겹친다: \emph{물리 가정 충돌}(평형 가역 이중웰 가정이 Si
지배 기작인 소성 경로의존 소산과 충돌 --- $55$ mV 상한이 수백 mV 를 못 담는 크기 불일치가 그 증상) $+$
\emph{수치해법 필요}(경로의존 구성식 $\sigma_h(\theta,\text{이력})$ 는 탄소성 모델로만 닫히며 이미
수치 모델\cite{koebbing2024,jiang_sihys2020}이 존재). 논리 공백 아님(식은 무결)$\cdot$정의상 implicit
아님(반전 문제 아님). 응력 구성식의 해석적 1차원리 접속은 본 장 범위 밖으로 선언한다 --- 유사 유도 날조
없이, 닫힌 데까지만.
\end{keybox}

이로써 생존 지도가 예고한 유일한 ``새 물리''(N3)의 취급이 정직하게 닫힌다: 응력 항의 \emph{선형 결합과
중심 이동}은 골격에 접속되나, 히스 gap 의 \emph{크기}는 소성 구성식(수치 소관)을 요구하는 물리 가정 충돌로
남는다. 이 진단은 Chapter 1 부록 예비 지도의 GS-1 선언(충돌 $+$ 미착수)을 한 걸음 진전시킨 것이다 ---
``미착수'' 를 ``부분 유도 닫힘 $+$ 구성식 수치 소관'' 으로 국소화했다.

% 자산(v1.0.22 R5 신규): [V22-R5-15 Larché–Cahn 선형 결합→평형 조건 결선 eq:si-lc-mu·eq:si-lc-shift(강체 중심이동 닫힘)]
%   [V22-R5-16 기계 히스 gap 구조=응력 경로차 eq:si-lc-gap(실측 −1.75 GPa 소성 앵커)]
%   [V22-R5-17 GS-1 공백 4분류 국소화 ssec:si-gs1(물리 가정 충돌+수치해법 필요·미착수→부분닫힘 진전)]
